\documentclass{article}
\usepackage{amsmath,amssymb}
\usepackage{enumitem}
\usepackage{listings}
\usepackage{xcolor}
\usepackage{hyperref}

\title{TreeLock Design Document}
\date{June 11, 2025}
\author{SC Dunn}

\definecolor{lightgray}{gray}{0.95}
\def\lstlanguagefiles{lstlean.tex}
\lstset{
  basicstyle=\ttfamily\small,
  backgroundcolor=\color{lightgray},
  frame=single,
  breaklines=true,
  postbreak=\mbox{\textcolor{red}{$\hookrightarrow$}\space}
}

\begin{document}

\maketitle

\section*{Project: TreeLock}
\textbf{Subtitle:} \textit{Merkle Tree-Based Inheritance in Decentralized Module Systems}

\section{Data Structure Design}

\subsection*{Tree Representation}
\begin{itemize}[noitemsep]
  \item \textbf{Decision:} Use a binary tree with \texttt{leaf} and \texttt{node} constructors.
  \item \textbf{Rationale:} Maps cleanly to Merkle semantics; enables recursive, verifiable hashing.
\end{itemize}

\begin{lstlisting}[language=lean]
inductive Tree (α : Type)
  | leaf : α → Tree α
  | node : Tree α → Tree α → Tree α
\end{lstlisting}

\section{Hashing Strategy}

\subsection*{Polymorphic Tree Hashing}
\begin{itemize}[noitemsep]
  \item \textbf{Decision:} Define \texttt{Hashable (Tree α)} given \texttt{Hashable α}.
  \item \textbf{Encoding:}
    \begin{itemize}
      \item \texttt{leaf x} $\to$ \texttt{hash [0x00, x]}
      \item \texttt{node l r} $\to$ \texttt{hash [0x01, hash l, hash r]}
    \end{itemize}
  \item \textbf{Rationale:} Ensures structure is encoded cryptographically.
\end{itemize}

\section{List to Tree Conversion}

\subsection*{Balanced Conversion}
\begin{itemize}[noitemsep]
  \item \textbf{Decision:} Split list in half recursively rather than flattening.
  \item \textbf{Rationale:} Guarantees balanced Merkle structure and efficient updates.
\end{itemize}

\begin{lstlisting}[language=Lean]
def listToTree : List Hash → Tree Hash
  | []      => Tree.leaf defaultHash
  | [x]     => Tree.leaf x
  | xs      =>
      let mid := xs.length / 2
      Tree.node (listToTree xs[..mid]) (listToTree xs[mid..])
\end{lstlisting}

\section{Merkle Hash Function}

\subsection*{Tree-Based Merkle Hash}
\begin{itemize}[noitemsep]
  \item \textbf{Decision:} Canonical definition via \texttt{Tree Hash}.
  \item \textbf{Rationale:} Simplifies termination and proof strategy.
\end{itemize}

\begin{lstlisting}[language=Lean]
def merkleHash (xs : List Hash) : Hash :=
  hash (listToTree xs)
\end{lstlisting}

\section{Formal Verification in Lean}

\subsection*{Theorems Proven}
\begin{itemize}[noitemsep]
  \item \texttt{merkleHash [] = defaultHash}
  \item \texttt{merkleHash [x] = hash [0x00, x]}
  \item \texttt{merkleHash [x, y] = hash [0x01, hash [0x00, x], hash [0x00, y]]}
\end{itemize}

\subsection*{Pending}
\begin{itemize}[noitemsep]
  \item Proof of consistency: \texttt{hash (listToTree xs) = merkleHash xs}
  \item General correctness of listToTree structure
\end{itemize}

\section{Module System Integration (Tangl)}

\subsection*{Override as Tree Update}
\begin{itemize}[noitemsep]
  \item \textbf{Decision:} Symbol overrides modeled as tree mutations.
  \item \textbf{Rationale:} Maintains integrity and traceability.
\end{itemize}

\subsection*{Recursive Lineage Tracking}
\begin{itemize}[noitemsep]
  \item \textbf{Decision:} TreeHash used as provenance signature.
  \item \textbf{Rationale:} Enables reproducibility and diffing.
\end{itemize}

\section{Hash Encoding Format}

\begin{itemize}[noitemsep]
  \item \textbf{Leaf tag:} \texttt{0x00}
  \item \textbf{Node tag:} \texttt{0x01}
  \item \textbf{Rationale:} Avoids collisions and domain confusion.
\end{itemize}

\section{Prototyping and Implementation}

\begin{itemize}[noitemsep]
  \item \textbf{Language:} Lean 4 (for executable proofs)
  \item \textbf{Prototype Frontend:} Tangl (mentioned only in prototype scope)
\end{itemize}

\section{Documentation and Paper Integration}

\begin{itemize}[noitemsep]
  \item \textbf{Appendix:} Appendix B describes formal connections to module systems.
  \item \textbf{Main Text:} Figures visualize tree hashing. Notation consistent across Lean + prose.
\end{itemize}

\section{Open Questions}

\begin{itemize}
  \item Generalization to non-binary or variadic branching trees
  \item Efficient diffing and reconciliation between tree versions
  \item Reusable tree patch or override languages
\end{itemize}

\end{document}
